% concatenated from CE20.08.tex
\documentclass[11pt]{amsart}

\usepackage[T1]{fontenc}
\usepackage[utf8]{inputenc}
\usepackage{amsmath,amssymb,amsfonts,mathtools,url}
\usepackage{geometry}
\usepackage{booktabs}
\usepackage{enumitem}
\usepackage{hyperref}
\hypersetup{hidelinks}
\usepackage[nameinlink,capitalize]{cleveref}
\geometry{margin=1in}

%\usepackage{refcheck}

\newcommand{\PP}{\mathbb P}
\newcommand{\CC}{\mathbb C}
\newcommand{\QQ}{\mathbb Q}
\newcommand{\NN}{\mathbb N}
\newcommand{\cA}{\mathcal A}
\newcommand{\cB}{\mathcal B}
\newcommand{\cM}{\mathcal M}
\newcommand{\AR}{\operatorname{AR}}
\newcommand{\mdr}{\operatorname{mdr}}
\newcommand{\rank}{\operatorname{rank}}
\newcommand{\Der}{\operatorname{Der}}
\newcommand{\Syz}{\operatorname{Syz}}
\newcommand{\Jac}{\operatorname{Jac}}
\newcommand{\indeg}{\operatorname{indeg}}
\newcommand{\HT}{\operatorname{HT}}
\newcommand{\hType}{\operatorname{hType}}
\newcommand{\HF}{\operatorname{HF}}

\theoremstyle{plain}
\newtheorem{theorem}{Theorem}[section]
\newtheorem{proposition}[theorem]{Proposition}
\newtheorem{lemma}[theorem]{Lemma}
\newtheorem{corollary}[theorem]{Corollary}
\newtheorem{conjecture}[theorem]{Conjecture}
\theoremstyle{definition}
\newtheorem{definition}[theorem]{Definition}
\newtheorem{example}[theorem]{Example}
\newtheorem{remark}[theorem]{Remark}

\numberwithin{equation}{section}

\title[A nine-line counterexample]{A nine-line counterexample to a conjecture on the minimal degree of Jacobian relations}

\author{Alexandru Dimca}
\address{Universit\'e C\^ote d'Azur, CNRS, LJAD, France  and Simion Stoilow Institute of Mathematics, P.O. Box 1-764\\ RO-014700 Bucharest, Romania}
\email{Alexandru.DIMCA@univ-cotedazur.fr}

\author{Piotr Pokora}
\address{Department of Mathematics, UKEN Krakow, Podchor\c a\.zych 2, PL-30-084 Krak\'ow, Poland}
\email{piotr.pokora@uken.krakow.pl}

\date{\today}

\begin{document}

\begin{abstract}
We construct two arrangements of nine lines in the complex projective plane with isomorphic intersection lattices but with different minimal degrees of Jacobian relations.  The common weak combinatorics is
\[
(n_2,n_3,n_4)=(9,7,1),
\]
so the example is far from the classical Ziegler--Yuzvinsky pair, whose weak combinatorics is \((n_{2},n_{3}) = (18,6)\).  For the two defining equations \(f\) and \(g\) we prove
\[
\mdr(f)=4,\qquad \mdr(g)=5.
\]
Since the degree is \(d=9\), the first equality gives \(\mdr(f)<d/2\).  Hence the pair gives a counterexample to a conjecture by the first author, that would have implied Terao's Conjecture about the combinatorial nature of the freeness property of line arrangements. Note that a very similar, but mildly weaker Conjecture, that is recalled below, also implies Terao's Conjecture and  is still open.
\end{abstract}

\maketitle

\section{Introduction}
Let \(S=\CC[x,y,z]\), and let \(C:f=0\) be a reduced plane curve of degree \(d\).  The Jacobian ideal is
\[
J_f=(f_x,f_y,f_z)\subset S,
\]
and the Milnor algebra is defined as
\[
M(f)=S/J_f.
\]
The graded module of Jacobian syzygies is
\[
\AR(f)=\{(a,b,c)\in S^3\;:\; a f_x+b f_y+c f_z=0\}.
\]
Equivalently, \(\AR(f)\) is the module \(D_0(f)\) of logarithmic derivations annihilating \(f\), after the usual identification
\[
(a,b,c)\longleftrightarrow a\partial_x+b\partial_y+c\partial_z,
\]
see \cite[Section 8.1]{DimcaBook}. The first numerical invariant of this module is the minimal degree of a Jacobian relation
\[
\mdr(f)=\min\{q\in\NN:\AR(f)_q\neq 0\}.
\]

For line arrangements, \(\mdr(f)\) is closely related to the freeness problem and to Terao-type questions.  Recall that the line arrangement $\cA$ is free if the graded $S$-module $D_0(f)$ is free. The celebrated Terao's Conjecture claims that if $\cA$ and $\cB$ are line arrangements in $\PP^2$ such that $\cA$ is free and $\cA$ and $\cB$ have isomorphic intersection lattices, then $\cB$ is also free. This conjecture has attracted a lot of interest, and it is known to hold for arrangements involving at most 14 lines, see \cite{BK, DIM13}.

If \(\cA:f=0\) is an arrangement of \(d\) lines in \(\PP^2\), then the intersection lattice determines the singularity multiplicities \(n_k\), hence also the total Tjurina number
\[
\tau(\cA)=\sum_{k\geq 2}n_k(k-1)^2.
\]
It does not, however, determine all Jacobian-syzygy data in general.  The classical examples of Ziegler \cite{Ziegler} and Yuzvinsky \cite{Yuzvinsky} give arrangements with the same intersection lattice but different Jacobian-syzygy behaviour; geometrically, the construction is related to a conic condition in Pascal-type configurations \cite{DimcaSticlaruPascal}.

The range \(\mdr(f)<d/2\) has been expected to be more rigid.  In particular, the first author formulated the following strengthening of Terao's conjecture for line arrangements.

\begin{conjecture}[{\cite[Conjecture~3.5]{DimcaOpen}}]
\label{mainc}
\emph{Let \(\cA:f=0\) be a line arrangement in \(\PP^2\), with \(d=\deg f\).  If \(\mdr(f)<d/2\), then \(\mdr(f)\) is combinatorially determined.}
\end{conjecture}

This conjecture would imply Terao's Conjecture, since the freeness of the arrangement $\cA$ is equivalent to 
$$\tau(\cA)=(d-1)^2-r(d-r-1),$$
where $r=\mdr(f)\leq (d-1)d/2$, see
\cite{duPCTC, Dmax}.

Note that the following weaker conjecture, but still implying Terao's Conjecture, is still open, see \cite[Conjecture~3.7]{DimcaOpen}.

\begin{conjecture}
\label{C2}
\emph{Let \(\cA:f=0\) be a line arrangement in \(\PP^2\), with \(d=\deg f\).  If \(\mdr(f)< (d-2)/2\), then \(\mdr(f)\) is combinatorially determined.}
\end{conjecture}
In view of the discussion at the beginning of Section 7, we see that Conjecture \ref{C2} is equivalent to claiming that the freeness defect $\nu(\cA)$ is determined by the combinatorics. In more geometrical terms,
Conjecture \ref{C2} is equivalent to claiming that the generic splitting type of the rank 2 vector bundle on $\PP^2$ associated to the gaded $S$-module $\AR(f)$ is determined by the combinatorics. This in turn is an open question asked in \cite[Question 7.12]{Cook+}, see for details \cite[Conjecture 3.9]{DimcaOpen}.

The purpose of the present note is to give an explicit counterexample to Conjecture \ref{mainc}. More precisely, we construct two arrangements of nine lines with isomorphic intersection lattices such that
\[
\mdr(f)=4,
\qquad
\mdr(g)=5.
\]
Since \(4<9/2\), the first arrangement satisfies the hypothesis of Conjecture \ref{mainc}, but the second arrangement has the same lattice and a different value of \(\mdr\).

The example is not a rephrasing of the classical Ziegler--Yuzvinsky pair.  That pair has only double and triple points, with weak combinatorics \((n_2,n_3)=(18,6)\).  Our arrangements have one quadruple point and seven triple points:
\[
(n_2,n_3,n_4)=(9,7,1).
\]
Thus the pair lies in the borderline nine-line case with maximal multiplicity \(4\).

A second goal of the paper is to record the homological nature of the two arrangements.  In the terminology of the hierarchy introduced in \cite{ADP}, the first arrangement has type \(1\), hence is plus-one generated \cite{AbePOG}, while the second has type \(2B\).  More concretely, their modules \(D_0(f)\) and \(D_0(g)\), and equivalently their Milnor algebras, have different minimal graded free resolutions.  Thus the counterexample is visible already at the level of \(\mdr\), and even more sharply at the level of homological type.

\section{Background on Jacobian syzygies and homological type}

We recall the conventions used throughout the paper.  Let \(C:f=0\) be a reduced plane curve of degree \(d\).  The module \(D_0(f)=\AR(f)\) is a graded \(S\)-module of rank two.  If its minimal homogeneous generators have degrees
\[
d_1\leq d_2\leq \cdots \leq d_m,
\]
we call this ordered sequence the sequence of exponents of \(C\).  In particular,
\[
d_1=\mdr(f).
\]
The minimal resolution of the Milnor algebra has the form
\begin{equation}\label{eq:Milnor-resolution-general}
0\longrightarrow F_3\longrightarrow F_2\longrightarrow S^3(-(d-1))\longrightarrow S\longrightarrow M(f)\longrightarrow 0,
\end{equation}
where the summands of \(F_2\) are obtained from the minimal generators of \(D_0(f)\).  More precisely, a generator of \(D_0(f)\) in degree \(e\) contributes a summand \(S(-(d-1+e))\) to \(F_2\).

Following \cite{ADP}, one defines the type of \(C\) by
\begin{equation}\label{eq:type-def}
t(C)=d_1+d_2+1-d.
\end{equation}
Equivalently, this number can be described as the initial degree of the Bourbaki ideal associated with a minimal generator of \(D_0(f)\).  Type \(0\) curves are precisely free curves, and type \(1\) curves are plus-one generated.  In type \(2\) there are two basic homological alternatives: type \(2A\), corresponding to a three-syzygy curve, and type \(2B\), corresponding to a four-syzygy curve.

For the purposes of this note, we shall use the following compact terminology.

\begin{definition}\label{def:homological-type}
Let \(\cA:f=0\) be a line arrangement.  Its \emph{homological type} is the collection
\[
\hType(\cA)=\bigl(d;\; \text{minimal resolution of }D_0(f);\; t(\cA)\bigr).
\]
When only the first part is needed, we write
\[
\HT_{D_0}(\cA)=\text{the minimal graded free resolution of }D_0(f).
\]
\end{definition}

This convention refines the numerical invariant \(\mdr(f)\).  Two arrangements may have the same intersection lattice and the same \(\mdr\), but still differ homologically through the graded Betti numbers of their Milnor algebras.  In the present example the distinction is stronger: the two arrangements already have different \(\mdr\)'s.

For later use, we also recall how \(\mdr\) is computed.  For each \(q\geq0\), consider the linear map
\begin{equation}\label{eq:Phi-map}
\Phi_{f,q}:S_q^3\longrightarrow S_{q+d-1},
\qquad
(a,b,c)\longmapsto a f_x+b f_y+c f_z.
\end{equation}
Then
\[
\AR(f)_q=\ker \Phi_{f,q}.
\]
Thus \(\mdr(f)\) is the smallest \(q\) for which \(\Phi_{f,q}\) fails to be injective.  Since
\[
\dim S_q=\binom{q+2}{2},
\]
this is an elementary exact linear-algebra computation over the ground field whenever the defining polynomial has rational coefficients.

\section{The two arrangements}

We work over \(\QQ\subset \CC\).  Let \(\cA:f=0\) be the arrangement of nine lines defined by
\begin{equation}\label{eq:f-def}
\begin{aligned}
f={}&(y-z)(y+2z)(2y+z)(x-2y-z)(x-y-2z)  \\
&\times (x-y+z)(x+y-z)(x+y+2z)(2x-y-2z).
\end{aligned}
\end{equation}
Let \(\cB:g=0\) be the arrangement defined by
\begin{equation}\label{eq:g-def}
\begin{aligned}
g={}&x(x-y)(x+y)(x-y-z)(y+z)z  \\
&\times (x-z)(x+y+2z)(x-2y-z).
\end{aligned}
\end{equation}
We label the lines in the displayed order.  Thus, for \(\cA\),
\[
\begin{array}{lll}
L_1:y-z=0, & L_2:y+2z=0, & L_3:2y+z=0,\\
L_4:x-2y-z=0, & L_5:x-y-2z=0, & L_6:x-y+z=0,\\
L_7:x+y-z=0, & L_8:x+y+2z=0, & L_9:2x-y-2z=0,
\end{array}
\]
and, for \(\cB\),
\[
\begin{array}{lll}
M_1:x=0, & M_2:x-y=0, & M_3:x+y=0,\\
M_4:x-y-z=0, & M_5:y+z=0, & M_6:z=0,\\
M_7:x-z=0, & M_8:x+y+2z=0, & M_9:x-2y-z=0.
\end{array}
\]

The following lemma verifies the combinatorics.

\begin{lemma}\label{lem:lattice}
With the labelling above, the non-double intersection points of both arrangements are exactly
\begin{equation}\label{eq:incidence-list}
\begin{gathered}
\{1,2,3\},\quad
\{1,4,5\},\quad
\{1,6,7\},\quad
\{2,4,6\},\\
\{3,5,7\},\quad
\{3,6,8\},\quad
\{4,7,9\},\quad
\{2,5,8,9\}.
\end{gathered}
\end{equation}
Consequently, \(\cA\) and \(\cB\) have isomorphic intersection lattices.
\end{lemma}

\begin{proof}
Represent a line \(a x+b y+c z=0\) by the vector \((a,b,c)\).  Three labelled lines \(i,j,k\) are concurrent if and only if the determinant of the corresponding \(3\times3\) matrix is zero.  The same determinant calculation for the two lists of line vectors gives exactly the subsets in \eqref{eq:incidence-list}.  The unique subset of cardinality four is \(\{2,5,8,9\}\), and every other dependent triple is one of the seven triples displayed above.  All other pairs meet in ordinary double points.

Thus the map \(L_i\mapsto M_i\) identifies the rank-two flats of the two arrangements and induces an isomorphism of intersection lattices.
\end{proof}

\begin{corollary}\label{cor:weak-combinatorics}
The common weak combinatorics of \(\cA\) and \(\cB\) is
\[
n_4=1,
\qquad
n_3=7,
\qquad
n_2=9.
\]
In particular,
\[
\tau(\cA)=\tau(\cB)=46.
\]
\end{corollary}


\section{The minimal degree of a Jacobian relation}

We now compute the kernels of the maps \(\Phi_{f,q}\) and \(\Phi_{g,q}\) from \eqref{eq:Phi-map}.  The computations are over \(\QQ\).  Since \(d=9\), the maps are
\[
\Phi_{h,q}:S_q^3\longrightarrow S_{q+8},
\qquad h\in\{f,g\}.
\]
The following table records the relevant ranks.

\begin{center}
\begin{tabular}{c|c|c|c|c|c}
\toprule
\(q\) & \(\dim S_q^3\) & \(\rank \Phi_{f,q}\) & \(\dim\AR(f)_q\) & \(\rank \Phi_{g,q}\) & \(\dim\AR(g)_q\) \\
\midrule
0 & 3  & 3  & 0 & 3  & 0 \\
1 & 9  & 9  & 0 & 9  & 0 \\
2 & 18 & 18 & 0 & 18 & 0 \\
3 & 30 & 30 & 0 & 30 & 0 \\
4 & 45 & 44 & 1 & 45 & 0 \\
5 & 63 & 59 & 4 & 59 & 4 \\
6 & 84 & 74 & 10 & 74 & 10 \\
7 & 108& 90 & 18 & 90 & 18 \\
\bottomrule
\end{tabular}
\end{center}

\begin{proposition}\label{prop:mdr}
For the arrangements \(\cA:f=0\) and \(\cB:g=0\) one has
\[
\mdr(f)=4,
\qquad
\mdr(g)=5.
\]
\end{proposition}

\begin{proof}
By definition, \(\AR(h)_q=\ker\Phi_{h,q}\).  The table shows that \(\Phi_{f,q}\) is injective for \(q=0,1,2,3\), while \(\dim\ker\Phi_{f,4}=1\).  Hence \(\mdr(f)=4\).  Similarly, \(\Phi_{g,q}\) is injective for \(q=0,1,2,3,4\), while \(\dim\ker\Phi_{g,5}=4\).  Hence \(\mdr(g)=5\).
\end{proof}

For completeness, we give one explicit degree-four syzygy for \(f\).  It has the form
\[
A f_x+B f_y+C f_z=0,
\]
where
\[
\begin{aligned}
A={}&-x(6x^3-17x^2y-10x^2z-14xy^2-20xyz-47xz^2 \\
&\hspace{2.2cm}+33y^3-15y^2z+75yz^2+69z^3),
\end{aligned}
\]
\[
\begin{aligned}
B={}&12x^3y-28x^2y^2+10x^2yz+18x^2z^2-4xy^3+20xy^2z \\
&-61xyz^2-36xz^3+12y^4-57y^3z+33y^2z^2+75yz^3+18z^4,
\end{aligned}
\]
and
\[
\begin{aligned}
C={}&12x^3z+18x^2y^2-x^2yz-17x^2z^2-40xy^2z-34xyz^2 \\
&-7xz^3-18y^4+39y^3z+69y^2z^2-3yz^3-6z^4.
\end{aligned}
\]

\section{Homological type}

The rank computations above also determine the minimal generators of the syzygy modules.  Let
\[
\mu_q(D_0(h))
=
\dim_\CC \frac{D_0(h)_q}{S_1D_0(h)_{q-1}}
\]
be the number of new minimal generators of \(D_0(h)\) in degree \(q\) modulo $S_1 D_{0}(h)_{q-1}$.  From the same linear algebra, one obtains
\[
\begin{array}{c|cccccccc}
q &0&1&2&3&4&5&6&7\\
\hline
\mu_q(D_0(f))&0&0&0&0&1&1&1&0\\
\mu_q(D_0(g))&0&0&0&0&0&4&0&0
\end{array}.
\]
Indeed, for \(f\), the unique degree-four syzygy contributes three independent multiples in degree five; since \(\dim D_0(f)_5=4\), there is exactly one new generator in degree five.  Similarly, comparing \(D_0(f)_6\) with \(S_1D_0(f)_5\) gives exactly one new generator in degree six.  No new generators occur in degree seven.  For \(g\), all four degree-five syzygies are new, and they generate the subsequent pieces.

\begin{proposition}\label{prop:D0-resolutions}
The minimal graded free resolutions of the two syzygy modules are
\begin{equation}\label{eq:D0-f-resolution}
0\longrightarrow S(-7)
\longrightarrow S(-4)\oplus S(-5)\oplus S(-6)
\longrightarrow D_0(f)
\longrightarrow 0,
\end{equation}
and
\begin{equation}\label{eq:D0-g-resolution}
0\longrightarrow S(-6)^2
\longrightarrow S(-5)^4
\longrightarrow D_0(g)
\longrightarrow 0.
\end{equation}
Consequently, \(\cA\) has type \(1\), with exponents \((4,5,6)\), while \(\cB\) has type \(2B\), with exponents \((5,5,5,5)\).
In particular,
\[
\hType(\cA)\neq \hType(\cB).
\]
\end{proposition}

\begin{proof}
The minimal generator degrees are read from the numbers \(\mu_q\) above.  Hence \(D_0(f)\) has minimal generators in degrees \(4,5,6\), while \(D_0(g)\) has four minimal generators, all in degree \(5\).

Since \(D_0(f)\) and \(D_0(g)\) have rank two, their Hilbert series determine the degrees of the first syzygies among these minimal generators.  For \(f\), the three generators in degrees \(4,5,6\) contribute
\[
\binom{7-4+2}{2}+\binom{7-5+2}{2}+\binom{7-6+2}{2}=10+6+3=19
\]
to degree \(7\), while the rank table gives \(\dim D_0(f)_7=18\).  Thus there is one relation in degree \(7\), giving \eqref{eq:D0-f-resolution}.  For \(g\), four degree-five generators contribute \(4\dim S_1=12\) to degree \(6\), while \(\dim D_0(g)_6=10\).  Thus there are two independent relations in degree \(6\), giving \eqref{eq:D0-g-resolution}.

Finally, using \eqref{eq:type-def}, for \(f\) we get
\[
t(\cA)=4+5+1-9=1,
\]
so \(\cA\) is plus-one generated.  For \(g\),
\[
t(\cB)=5+5+1-9=2.
\]
Since \(D_0(g)\) has four minimal generators and the two first relations have the expected linear form in the type-two case, this is type \(2B\) in the terminology of \cite{ADP}.
\end{proof}
\begin{remark}
In the light of \cite[Definition 1.1]{mpog}, the arrangement $\mathcal{A}$ is minimal plus-one generated, which means that $d_{1}+d_{2} = {\rm deg}(f)$ and $d_{3}=d_{2}+1$.
\end{remark}
Shifting \eqref{eq:D0-f-resolution} and \eqref{eq:D0-g-resolution} by \(d-1=8\), we obtain the corresponding resolutions of the Milnor algebras.

\begin{corollary}\label{cor:Milnor-resolutions}
The Milnor algebra \(M(f)=S/J_f\) has minimal graded free resolution
\begin{equation}\label{eq:Mf-resolution}
0\longrightarrow S(-15)
\longrightarrow S(-12)\oplus S(-13)\oplus S(-14)
\longrightarrow S(-8)^3
\longrightarrow S
\longrightarrow M(f)
\longrightarrow 0.
\end{equation}
The Milnor algebra \(M(g)=S/J_g\) has minimal graded free resolution
\begin{equation}\label{eq:Mg-resolution}
0\longrightarrow S(-14)^2
\longrightarrow S(-13)^4
\longrightarrow S(-8)^3
\longrightarrow S
\longrightarrow M(g)
\longrightarrow 0.
\end{equation}
\end{corollary}

\begin{remark}
The Hilbert functions of the two Milnor algebras are also different.  From \eqref{eq:Mf-resolution} and \eqref{eq:Mg-resolution}, the Hilbert series are
\[
H_{M(f)}(t)=\frac{1-3t^8+t^{12}+t^{13}+t^{14}-t^{15}}{(1-t)^3},
\]
and
\[
H_{M(g)}(t)=\frac{1-3t^8+4t^{13}-2t^{14}}{(1-t)^3}.
\]
For instance,
\[
\dim M(f)_{12}=47,
\qquad
\dim M(g)_{12}=46.
\]
Both Hilbert functions stabilize to the common Tjurina number \(46\).
\end{remark}

\section{Counterexample to Conjecture~3.5}

We now state the conclusion explicitly.

\begin{theorem}\label{thm:counterexample}
The pair \((\cA,\cB)\) defined by \eqref{eq:f-def} and \eqref{eq:g-def} is a pair of nine-line arrangements with isomorphic intersection lattices but with
\[
\mdr(\cA)=4,
\qquad
\mdr(\cB)=5.
\]
Consequently, Conjecture~\ref{mainc} is false.
\end{theorem}

\begin{proof}
By Lemma \ref{lem:lattice}, the arrangements have isomorphic intersection lattices.  By Proposition \ref{prop:mdr}, their minimal degrees of Jacobian relations are \(4\) and \(5\), respectively.  Since the degree is \(d=9\), the first arrangement satisfies
\[
\mdr(\cA)=4<\frac92=\frac d2.
\]
Conjecture \ref{mainc} predicts that, under this strict inequality, \(\mdr\) is determined by the intersection lattice.  But \(\cB\) has the same intersection lattice as \(\cA\) and a different value of \(\mdr\).  This contradicts the conjecture.
\end{proof}

\section{The freeness defect \texorpdfstring{\(\nu\)}{nu}}

It is useful to separate the behaviour of \(\mdr\) from the behaviour of the freeness defect \(\nu\).  By the formulas recalled in \cite{DimcaOpen}, if \(C:f=0\) has degree \(d\), total Tjurina number \(\tau(C)\), and \(r=\mdr(f)\), then
\[
\nu(C)=(d-1)^2-r(d-1-r)-\tau(C)
\]
when \(r<(d-2)/2\), while for  \(r\geq (d-2)/2\) one has
\[
\nu(C)=\left\lceil \frac34(d-1)^2\right\rceil-\tau(C).
\]
In the present example \(d=9\) and \(\tau=46\).  We get

\[
\nu(\cA)=\nu(\cB)=\left\lceil \frac34\cdot 8^2\right\rceil-46=48-46=2.
\]
The example therefore attacks precisely the combinatoriality of \(\mdr\) in the range of Conjecture \ref{mainc}.  It does not, by itself, disprove the corresponding combinatoriality statement for \(\nu\)  in Conjecture \ref{C2}.
\begin{remark}
The pair \((\mathcal A,\mathcal B)\) realizes the only two possibilities for curves
with \(\nu=2\) described in \cite[Remark~1.6]{mpog}: the arrangement
\(\mathcal A\) is minimal plus-one generated, whereas \(\mathcal B\) is a
maximal Tjurina curve of type \((2r-1,r)\).
\end{remark}

\section*{Funding}
Alexandru Dimca is partially supported from the project ``Singularities and Applications'' - CF 132/31.07.2023 funded by the European Union - NextGenerationEU - through Romania's National Recovery and Resilience Plan.

Piotr Pokora is supported by the National Science Centre (Poland) Sonata Bis Grant 
\[\textbf{2023/50/E/ST1/00025.} \] For the purpose of Open Access, the author has applied a CC-BY public copyright license to any Author Accepted Manuscript (AAM) version arising from this submission.

\begin{thebibliography}{99}

\bibitem{AbePOG}
T. Abe,
\emph{Plus-one generated and next to free arrangements of hyperplanes},
Int. Math. Res. Not. IMRN \textbf{2021}, no.~12, 9233--9261.

\bibitem{ADP}
T. Abe, A. Dimca and P. Pokora,
\emph{A new hierarchy for complex plane curves},
Canad. Math. Bull., published online 2025, 1--25. \texttt{doi:10.4153/S0008439525101422}

\bibitem{BK}  M. Barakat, L. K\" uhne, Computing the nonfree locus of the moduli space of arrangements and Terao’s freeness conjecture.
Math. Comp. 92: 1431--1452 (2023).

\bibitem{Cook+} D. Cook, B. Harbourne, J. Migliore, U. Nagel,  Line arrangements and configurations of points with an unexpected geometric property. {Compositio Math.} {154}(2018), 2150--2194.

\bibitem{DimcaBook}
A. Dimca,
\emph{Hyperplane arrangements: an introduction},
Universitext, Springer, Cham, 2017.

\bibitem{Dmax}  A. Dimca, Freeness versus maximal global Tjurina number for plane curves, 
Math. Proc. Cambridge Phil. Soc.  163 (2017), 161--172.

\bibitem{DimcaOpen}
A. Dimca,
\emph{On free curves and related open problems},
Rev. Roumaine Math. Pures Appl. \textbf{69} (2024), no.~2, 129--150.

\bibitem{DIM13} A. Dimca, D. Ibadula, A. M\u acinic,
Freeness for 13 lines arrangements is combinatorial, Discrete Mathematics 342 (2019), 2445--2453.

\bibitem{DimcaSticlaruPascal}
A. Dimca and G. Sticlaru,
\emph{From Pascal's theorem to the geometry of Ziegler's line arrangements},
J. Algebraic Combin. \textbf{60} (2024), 991--1009.

\bibitem{mpog}
A. Dimca and G. Sticlaru, 
\textit{Plus-one generated curves, Briançon-type polynomials and eigenscheme ideals}, Result. Math. \textbf{80(2)} (2025), Paper No. 51, 22 p.

\bibitem{duPCTC} A.A. du Plessis,  C.T.C. Wall, Application of the theory of the
discriminant to highly singular plane curves, Math. Proc. Cambridge
Phil. Soc.,  126(1999), 259-266. 

\bibitem{Yuzvinsky}
S. Yuzvinsky,
\emph{The first two obstructions to the freeness of arrangements},
Trans. Amer. Math. Soc. \textbf{335} (1993), no.~1, 231--244.

\bibitem{Ziegler}
G. M. Ziegler,
\emph{Combinatorial construction of logarithmic differential forms},
Adv. Math. \textbf{76} (1989), 116--154.

\end{thebibliography}

\end{document}


\appendix

\section{SINGULAR verification}

The following SINGULAR script verifies the incidence data, the dimensions of the graded syzygy spaces, and the minimal resolutions.  The rank computation is deliberately written in a direct linear-algebra form; this makes the verification of \(\mdr\) independent of any specialized package.

\begin{verbatim}
///////////////////////////////////////////////////////////////
// Nine-line counterexample to combinatoriality of mdr
///////////////////////////////////////////////////////////////

ring R = 0,(x,y,z,u(1..84)),dp;

poly f =
 (y-z)
*(y+2*z)
*(2*y+z)
*(x-2*y-z)
*(x-y-2*z)
*(x-y+z)
*(x+y-z)
*(x+y+2*z)
*(2*x-y-2*z);

poly g =
 x
*(x-y)
*(x+y)
*(x-y-z)
*(y+z)
*z
*(x-z)
*(x+y+2*z)
*(x-2*y-z);

proc fac(int n)
{
    int ans = 1;
    int i;
    for (i=2; i<=n; i=i+1) { ans = ans*i; }
    return(ans);
}

proc nmons(int d)
{
    return((d+1)*(d+2)/2);
}

proc genericForm(int d, int offset)
{
    poly F = 0;
    int i; int j; int k; int cnt = 0;
    for (i=0; i<=d; i=i+1)
    {
        for (j=0; j<=d-i; j=j+1)
        {
            k = d-i-j;
            cnt = cnt+1;
            F = F + var(offset+cnt)*x^i*y^j*z^k;
        }
    }
    return(F);
}

proc coeffMon(poly P, int i, int j, int k)
{
    poly Q = P;
    int a;
    int den;
    for (a=1; a<=i; a=a+1) { Q = diff(Q,x); }
    for (a=1; a<=j; a=a+1) { Q = diff(Q,y); }
    for (a=1; a<=k; a=a+1) { Q = diff(Q,z); }
    Q = subst(Q,x,0);
    Q = subst(Q,y,0);
    Q = subst(Q,z,0);
    den = fac(i)*fac(j)*fac(k);
    Q = Q/den;
    return(Q);
}

proc coeffIdeal(poly P, int d)
{
    ideal E;
    poly c;
    int i; int j; int k; int cnt = 0;
    for (i=0; i<=d; i=i+1)
    {
        for (j=0; j<=d-i; j=j+1)
        {
            k = d-i-j;
            c = coeffMon(P,i,j,k);
            if (c <> 0)
            {
                cnt = cnt+1;
                E[cnt] = c;
            }
        }
    }
    return(E);
}

proc ARdim(poly h, int r)
{
    poly A = genericForm(r,3);
    poly B = genericForm(r,31);
    poly C = genericForm(r,59);
    poly P = A*diff(h,x) + B*diff(h,y) + C*diff(h,z);
    ideal E = coeffIdeal(P,r+8);
    ideal G = std(E);
    int N = nmons(r);
    int rk = size(G);
    int dim = 3*N - rk;
    return(dim);
}

proc printARdims(poly h, int maxr)
{
    int r;
    for (r=0; r<=maxr; r=r+1)
    {
        print("degree " + string(r) +
              " : dim AR = " + string(ARdim(h,r)));
    }
}

print("Arrangement A");
printARdims(f,7);

print("Arrangement B");
printARdims(g,7);

///////////////////////////////////////////////////////////////
// Expected output:
//
// Arrangement A:
// degree 0,1,2,3 : 0
// degree 4       : 1
// degree 5       : 4
// degree 6       : 10
// degree 7       : 18
// Hence mdr(f)=4.
//
// Arrangement B:
// degree 0,1,2,3,4 : 0
// degree 5         : 4
// degree 6         : 10
// degree 7         : 18
// Hence mdr(g)=5.
///////////////////////////////////////////////////////////////

// For minimal resolutions one may work in the smaller ring:
ring T = 0,(x,y,z),dp;
poly ff = imap(R,f);
poly gg = imap(R,g);

ideal Jf = jacob(ff);
ideal Jg = jacob(gg);

module ARf = syz(Jf);
module ARg = syz(Jg);

print("Betti table for D0(f):");
betti(minres(res(ARf,0)));

print("Betti table for D0(g):");
betti(minres(res(ARg,0)));

print("Betti table for M(f)=S/Jf:");
betti(minres(res(std(Jf),0)));

print("Betti table for M(g)=S/Jg:");
betti(minres(res(std(Jg),0)));
\end{verbatim}
